\clearpage
\SourceLeaf{70}
\pdfbookmark[1]{Observationes cometae secundi a. 1813 — Anzeige 1814}{anzeige-cometae1814}
\begin{center}\rule{.75\textwidth}{.3pt}\\
Göttingische gelehrte Anzeigen. 1814. Januar 3.\\
\rule{.75\textwidth}{.3pt}\end{center}
Wir haben noch die Anzeige einer kleinen vom Prof. \textsc{Gauss} am 10. Sept.
des vorigen Jahrs der königl. Societät überreichten Abhandlung nachzuholen,
überschrieben:
\begin{center}
\textit{Observationes cometae secundi a.} 1813 \textit{in observatorio Gottingensi factae, adiectis
nonnullis adnotationibus circa calculum orbitarum parabolicarum.}
\end{center}
Dieser Comet wurde bekanntlich am 3. April des vorigen Jahrs von unserm
Herrn Prof. \textsc{Harding} im Poniatovskischen Stier entdeckt: die auf der hiesigen
Sternwarte angestellten Beobachtungen desselben gehen vom 9. bis 25. April.
Man findet sie in der Abhandlung \textit{vollständig}: da indessen der grössere Theil derselben
auch schon an einem andern Ort bekannt gemacht ist, so übergehen wir
sie hier mit Stillschweigen, und heben aus der Abhandlung nur die parabolische
Bahn aus, welche Prof. \textsc{Gauss} den Herrn Dr. \textsc{Gerling} in Cassel nach den hiesigen
und einigen Bremer und Pariser Beobachtungen zu berechnen veranlasste:
\begin{center}\small
Durchgangszeit durch die Sonnennähe 1813 Mai 19. $10^{\mathrm h}\;24^{\mathrm m}\;5^{\mathrm s}$\\
M. Z. in Göttingen
\end{center}
\begin{center}\small
\setlength{\tabcolsep}{3pt}
\begin{tabular}{ll}
Länge der Sonnennähe $\quad.\;.\;.\;.$ & $197^0\;43'\;7''7$\\
Länge des aufsteigenden Knoten $\quad,$ & $42\;40\;15.2$\\
Neigung der Bahn $\quad.\;.\;.\;.\;.\;.$ & $81\quad2\;11.8$\\
Logarithm des kleinsten Abstandes & $0.0849212$\\
\multicolumn{2}{l}{Bewegung \textit{rückläufig.}}
\end{tabular}\end{center}

\SourceLeaf{71}%
Noch vor wenigen Jahrzehnden war die Anzahl der Personen in ganz Europa,
die eine Cometenbahn zu berechnen im Stande waren, nur klein: gegenwärtig
ist dieses Geschäft durch vervollkommnete Methoden so erleichtert und
vereinfacht, dass ein sonst fähiger Kopf sich ohne Schwierigkeit damit vertraut
machen, und in weniger Stunden, als sonst Tage erforderlich waren, eine Cometenbahn
bestimmen kann. Zur \textit{ersten} Berechnung einer Bahn aus drei Beobachtungen
lässt die bequeme Methode des Herrn Dr. \textsc{Olbers} fast nichts zu wünschen
übrig, wenn man anders Beobachtungen wählen kann, bei welchen die
Richtung der geocentrischen Bewegung nicht zu nahe an die Richtung des grössten
Kreises fällt, welcher durch den mittelsten geocentrischen Ort des Cometen
und den entsprechenden Sonnenort auf der Himmelskugel gezogen wird. Der
zweite Theil der vorliegenden Abhandlung hat zum Zweck, einige Punkte dieser
schönen Methode noch etwas mehr zu vereinfachen, und zur numerischen
Berechnung noch etwas bequemer zu machen. Dieser Versuch betrifft das indirecte
Verfahren, wodurch genäherte Werthe der Abstände des Cometen von der
Sonne und der Chorde zwischen den beiden äussern Örtern bestimmt werden.
Dass diess Verfahren indirect ist, könnte nur ein im astronomischen Calcul Unerfahrner
der Methode zum Vorwurf machen: aber die Einrichtung, welche ihm
\textsc{Olbers} gegeben hat, ist mit der Unbequemlichkeit verbunden, dass in den meisten
Fällen, auch wenn man übrigens, wie billig bei einer ersten Näherung, nicht
die äusserste Schärfe in die Rechnung legen will, doch die kleinern Tafeln mit
fünf Decimalen dabei zu grosse Fehler hervorbringen würden, weil mehrere ihrer
Natur nach sehr kleine Grössen in der Form von Differenzen anderer, welche
beinahe gleich werden, erscheinen. Prof. \textsc{Gauss} suchte dieser Unbequemlichkeit
dadurch zu begegnen, dass er durch Einführung mehrerer Hülfsgrössen,
welche vermittelst der kleinen Tafeln immer hinreichend genau bestimmt werden
können, den Formeln für die Abstände von der Sonne und die Chorde eine andere
Gestalt gab, bei welcher auch die nachmaligen Versuche der indirecten Auflösung
noch etwas an Bequemlichkeit zu gewinnen scheinen. Einen Auszug vertragen
diese Umformungen nicht, wenn wir nicht diesen Theil der Abhandlung
hier ganz abschreiben wollen: wir verweisen also auf den nächstens erscheinenden
zweiten Band der Commentationen der Societät, in welchem die ganze Abhandlung
bereits abgedruckt ist.

Der Verfasser glaubte manchem einen Dienst zu erweisen, wenn er mit den
